\chapter{Introdução}
A modelagem de sistemas complexos fornece o arcabouço matemático e computacional necessário para entender como interações locais entre componentes individuais geram comportamentos globais emergentes. Recentemente, a adoção de Modelos de Linguagem de Grande Escala (LLMs) impulsionou a criação de sistemas multiagentes autônomos. No entanto, o estudo da dinâmica de rede formada por múltiplos agentes que operam de maneira descentralizada (edge computing) e baseada em LLMs de código aberto executados localmente ainda é incipiente.

\section{Justificativa}
A dependência de APIs centralizadas de modelos de linguagem cria gargalos de privacidade, latência e custo. Ao investigar redes de agentes rodando inferência localmente, este projeto contribui para a literatura de sistemas complexos aplicados à inteligência artificial distribuída, unindo o rigor formal matemático às necessidades reais de infraestruturas tecnológicas modernas.

\subsection{Objetivo Geral}
Modelar e simular as interações em uma rede complexa de agentes autônomos impulsionados por LLMs quantizados (como Dolphin-Mistral e Qwen), analisando a emergência de fenômenos como consenso, polarização e difusão de informação sob restrições de hardware e conectividade local.

\subsection{Objetivos Específicos}
\begin{itemize}
    \item Desenvolver um modelo matemático baseado em grafos e álgebra linear para representar o fluxo de informações probabilísticas entre os agentes.
    \item Implementar um ambiente de simulação em Python e Rust para orquestrar a comunicação descentralizada entre instâncias de LLMs locais.
    \item Analisar o impacto da topologia da rede nas taxas de convergência de ideias ou resolução cooperativa de problemas complexos.
\end{itemize}
